\section{Metodologia}

O experimento visou estimar a capacitância real do capacitor em um Circuito RC série e determinar a impedância equivalente de um Circuito RLC misto. Para o circuito RC, utilizou-se uma \textit{protoboard} para montar o sistema, um resistor $R_{ext}=100~\Omega$ e um gerador de funções com impedância interna $R_{int}=50~\Omega$, resultando em uma resistência efetiva $R_{ef}=150~\Omega$, além de um capacitor com capacitância nominal de \(C = \qty{330}{nF} \). A constante de tempo de relaxação $(\tau_{exp}=124~\mu s)$ foi medida com um osciloscópio, verificando o tempo para a tensão no capacitor decair a $1/e$ de seu valor inicial. A capacitância real foi calculada pela relação $\tau=R\cdot C$. O circuito RLC misto consistia em um resistor (\(R = \qty{100}{\Omega} \)) em série com um capacitor (\(C = \qty{3.3}{\micro F}\)) e um indutor (\(L = \qty{15}{\micro H} \)) em paralelo entre si.

